\chapter{Metodologia di sviluppo e tecnologie adottate}

\section{Metodologia di sviluppo}

Lo sviluppo del sistema \textit{JarFin} è stato condotto adottando una metodologia Agile, con l’obiettivo di favorire un processo di progettazione incrementale, flessibile e orientato al cambiamento. In particolare, il progetto ha seguito i principi di \textbf{SCRUM}, integrati con l’approccio \textbf{Agile Model Driven Development (AMDD)}, al fine di bilanciare la necessità di una progettazione architetturale solida con i vantaggi dell’iterazione continua e del feedback costante.

SCRUM è stato utilizzato come framework di riferimento per l’organizzazione del lavoro, la pianificazione delle attività e la gestione delle iterazioni di sviluppo. Il ciclo di vita del progetto è stato articolato in sprint di durata indicativa pari a due settimane, ciascuno caratterizzato da obiettivi ben definiti e dalla produzione di un incremento funzionale del sistema. Questa impostazione ha consentito di monitorare costantemente l’avanzamento del progetto e di adattare le scelte progettuali in base alle esigenze emerse durante lo sviluppo \cite{scrumOrg}.

\begin{figure}[H]
    \centering
    \includegraphics[width=0.9\textwidth]{2.0_scrum_framework.png}
    \caption{Il framework SCRUM e il ciclo degli Sprint.}
    \label{fig:scrum_framework}
\end{figure}

L’approccio AMDD ha guidato le attività di modellazione e progettazione. Invece di definire in modo esaustivo l’intera architettura nelle fasi iniziali, la modellazione è stata condotta in maniera incrementale, producendo modelli essenziali e sufficienti a supportare le decisioni architetturali di ciascuna iterazione. In particolare, sono stati realizzati diagrammi UML per i casi d’uso più rilevanti, successivamente raffinati nel corso dello sviluppo. Questo approccio ha permesso di evitare un’eccessiva rigidità progettuale, mantenendo una visione globale del sistema.

\section{Organizzazione delle iterazioni}

Il ciclo di sviluppo di JarFin è stato articolato in una serie di iterazioni successive, ciascuna focalizzata su specifici obiettivi funzionali e architetturali:

\begin{itemize}
    \item \textbf{Iterazione 0 – Envisioning e Analisi}: definizione del dominio applicativo, raccolta e analisi dei requisiti funzionali e non funzionali, identificazione dei principali casi d’uso e progettazione dell’architettura di alto livello;
    \item \textbf{Iterazione 1 – Core Services}: sviluppo del servizio di accounting e implementazione della persistenza dei dati.
    \item \textbf{Iterazione 2 – Analytics e Algoritmi}: sviluppo del servizio di analytics e implementazione degli algoritmi per l’elaborazione e l’analisi dei dati finanziari;
    \item \textbf{Iterazione 3 – Integrazione e Interazione}: integrazione del modulo di Natural Language Understanding, dell’API Gateway e dell’interfaccia web.
\end{itemize}

Questa suddivisione ha consentito di affrontare progressivamente la complessità del sistema, validando le scelte progettuali a ogni iterazione e riducendo il rischio di errori architetturali nelle fasi avanzate dello sviluppo.

\section{Architettura generale del sistema}

L’architettura di JarFin è basata su un paradigma a microservizi, in cui il sistema è scomposto in un insieme di componenti indipendenti, ciascuno responsabile di una specifica area funzionale. Tale scelta architetturale è motivata dall’esigenza di garantire modularità, scalabilità e manutenibilità, oltre a facilitare l’evoluzione futura del sistema.

Ogni microservizio è progettato come un’unità autonoma, dotata di una propria logica applicativa e, ove necessario, di un proprio livello di persistenza. La comunicazione tra i servizi avviene tramite API REST e scambio di messaggi in formato JSON, secondo i principi dello stile architetturale REST \cite{fieldingRest}.

L’adozione di un’architettura distribuita consente inoltre di isolare le responsabilità dei singoli componenti, riducendo l’impatto delle modifiche e migliorando la resilienza complessiva del sistema in caso di malfunzionamenti localizzati.

\section{Tecnologie adottate}

La scelta dello stack tecnologico è stata guidata dalla necessità di utilizzare strumenti robusti, ampiamente supportati dalla community e adatti allo sviluppo di sistemi distribuiti.

\subsection{Spring Boot}

Per l’implementazione dei microservizi è stato utilizzato il framework \textbf{Spring Boot}, che fornisce un ambiente di sviluppo semplificato per la realizzazione di applicazioni Java basate su architettura a servizi. Spring Boot consente di ridurre significativamente la configurazione manuale, favorendo lo sviluppo rapido di applicazioni stand-alone e facilmente deployabili \cite{springboot}.

Il framework supporta nativamente la creazione di API REST, l’integrazione con sistemi di persistenza e la gestione della configurazione dei servizi, risultando particolarmente adatto allo sviluppo di architetture a microservizi.

\subsection{Database PostgreSQL e Containerizzazione}

La persistenza dei dati è affidata a un database relazionale \textbf{PostgreSQL}, scelto per la sua affidabilità, conformità agli standard SQL e supporto a funzionalità avanzate di gestione dei dati. PostgreSQL garantisce l’integrità delle transazioni finanziarie (proprietà ACID) e supporta operazioni di interrogazione efficienti, fondamentali per le funzionalità di analisi offerte dal sistema.

Per semplificare la gestione dell'ambiente di sviluppo e garantire l'isolamento dei dati, l'istanza del database è stata containerizzata tramite \textbf{Docker}.
Questa scelta architetturale consente di avviare il servizio di persistenza in modo rapido e indipendente dalla configurazione del sistema operativo ospite, evitando installazioni locali invasive e garantendo un ambiente sempre pulito e riproducibile (Figura~\ref{fig:docker_vm}).

\begin{figure}[H]
    \centering
    \includegraphics[width=0.8\textwidth]{2.1_docker_vs_VM.png}
    \caption{Confronto architetturale tra Macchine Virtuali e Container Docker \cite{geeksforgeeksDocker}.}
    \label{fig:docker_vm}
\end{figure}

\subsection{Modulo di Natural Language Understanding}

Per l’elaborazione del linguaggio naturale è stato sviluppato un modulo di Natural Language Understanding implementato in Java come microservizio dedicato. Il sistema utilizza tecniche di analisi lessicale e \textit{pattern matching} per identificare intenti e parametri rilevanti a partire dalle frasi dell’utente, consentendo l’esecuzione di operazioni finanziarie tramite comandi in linguaggio naturale.

La scelta di una soluzione interna, anziché basata su servizi cloud esterni, consente di ridurre la latenza, mantenere il controllo sui dati sensibili dell’utente e garantire una maggiore integrazione con l’architettura complessiva del sistema.

\section{Toolchain di progetto}

Lo sviluppo del sistema \textit{JarFin} è stato supportato da una toolchain articolata, composta da strumenti dedicati alla modellazione, allo sviluppo, al versioning, al testing e all’analisi della qualità del codice.
La scelta degli strumenti è stata guidata dall’obiettivo di garantire un processo di sviluppo strutturato, tracciabile e conforme alle buone pratiche dell’ingegneria del software.

\subsection{Strumenti di modellazione}

Per la modellazione del sistema è stato utilizzato \textbf{UMLet}, uno strumento leggero per la creazione di diagrammi UML. UMLet è stato impiegato per la realizzazione dei principali artefatti di modellazione previsti dall’approccio AMDD, in particolare:
\begin{itemize}
    \item diagrammi dei casi d’uso, per rappresentare le funzionalità offerte dal sistema e le interazioni con gli attori;
    \item diagrammi delle classi, per descrivere la struttura statica dei principali componenti software;
    \item diagrammi dei componenti, per modellare l’architettura a microservizi;
    \item diagrammi di sequenza, per analizzare il comportamento dinamico del sistema nei casi d’uso critici.
    \item diagrammi di deployment, per descrivere la distribuzione fisica dei nodi e l'architettura containerizzata del database.
\end{itemize}

L’utilizzo di UMLet ha consentito di produrre modelli essenziali e facilmente modificabili, in linea con i principi della modellazione incrementale promossi dall’Agile Model Driven Development.

\subsection{Versioning e Project Management}

Il controllo di versione del codice sorgente e la gestione del progetto sono stati affidati a \textbf{GitHub}.
Git è stato utilizzato come sistema di versioning distribuito, permettendo di tracciare l’evoluzione del codice, gestire rami di sviluppo separati e mantenere uno storico completo delle modifiche.

\newpage

GitHub è stato inoltre impiegato come strumento di supporto al project management, sfruttando:
\begin{itemize}
    \item repository centralizzato per il codice sorgente;
    \item Kanban board per la pianificazione delle attività e il tracciamento delle iterazioni;
    \item Pull Request per la revisione del codice e l’integrazione controllata delle funzionalità.
\end{itemize}

Questa impostazione ha favorito un processo di sviluppo ordinato e trasparente, coerente con la metodologia SCRUM adottata.

\subsection{Sviluppo backend}

Il backend del sistema è stato sviluppato utilizzando il linguaggio \textbf{Java} e il framework \textbf{Spring Boot}.
Spring Boot è stato scelto per la sua capacità di semplificare lo sviluppo di applicazioni basate su architettura a microservizi, riducendo la configurazione manuale e favorendo la realizzazione di servizi REST stand-alone.

Ogni microservizio di JarFin è implementato come applicazione Spring Boot indipendente, responsabile di una specifica area funzionale del sistema, in linea con i principi di separazione delle responsabilità e modularità.

\subsection{Testing delle API REST}

Per la verifica manuale e il testing degli endpoint REST è stato utilizzato \textbf{Postman}.
Lo strumento ha permesso di:
\begin{itemize}
    \item testare le operazioni CRUD esposte dai microservizi;
    \item validare il corretto formato delle richieste e delle risposte JSON;
    \item simulare scenari di utilizzo reali durante le fasi di sviluppo e integrazione.
\end{itemize}

L’uso di Postman ha facilitato l’individuazione precoce di errori di comunicazione tra i servizi e ha supportato la fase di debugging dell’API Gateway.

\subsection{Unit Testing e Mocking}

Il testing automatico della logica applicativa è stato realizzato tramite \textbf{JUnit 5}, framework di riferimento per il testing in ambiente Java, affiancato da \textbf{Mockito}, un framework di mocking essenziale per isolare le unità di codice.
Mentre JUnit è stato utilizzato per definire il ciclo di vita dei test e le asserzioni, Mockito ha permesso di simulare il comportamento delle dipendenze esterne (mock), garantendo che i test unitari si focalizzassero esclusivamente sulla logica interna dei servizi, in particolare nei moduli di accounting, analytics e NLU.

L’adozione combinata di questi strumenti ha contribuito a migliorare l’affidabilità del sistema e a ridurre il rischio di regressioni durante l’evoluzione del codice.

\subsection{Analisi della copertura del codice}

Per l’analisi della copertura dei test è stato utilizzato il plugin \textbf{EclEmma} (basato sulla libreria \textbf{JaCoCo}), integrato direttamente nell’ambiente di sviluppo.
Lo strumento consente di misurare dinamicamente la percentuale di codice eseguito durante l’esecuzione dei test (code coverage), fornendo indicatori visivi e quantitativi utili per valutare l’efficacia della suite di test.

Le metriche di copertura hanno supportato l’identificazione delle porzioni di codice critiche non ancora verificate, guidando il miglioramento progressivo dei test unitari.

\subsection{Analisi statica e architetturale}

La qualità del software è stata monitorata su due livelli distinti mediante l'uso combinato di strumenti di analisi statica:

\subsubsection{Qualità del Codice (SonarLint)}
La qualità strutturale e la conformità agli standard sono state analizzate mediante \textbf{SonarLint}. Lo strumento è stato utilizzato per:
\begin{itemize}
    \item rilevare metriche di complessità cognitiva e ciclomatica puntuale;
    \item individuare \textit{code smells}, bug potenziali e vulnerabilità di sicurezza;
    \item supportare la verifica del rispetto dei principi di Clean Code.
\end{itemize}

\subsubsection{Qualità Architetturale (CodeMR)}
Per un'analisi di più alto livello, è stato introdotto \textbf{CodeMR}, uno strumento di analisi statica avanzata focalizzato sulla struttura del software.
A differenza dei linter classici, CodeMR ha permesso di visualizzare e mappare le dipendenze tra i package, calcolando metriche architetturali fondamentali quali:
\begin{itemize}
    \item \textbf{Accoppiamento (Coupling)}: per valutare l'indipendenza dei moduli;
    \item \textbf{Complessità Ciclomatica}: aggregata a livello di classe e package;
    \item \textbf{Coesione}: per verificare la responsabilità singola delle classi.
\end{itemize}
L'utilizzo di CodeMR è stato determinante per verificare l'assenza di dipendenze cicliche tra i microservizi e per identificare preventivamente potenziali ``God Classes'' (classi monolitiche con troppe responsabilità), facilitando il refactoring verso un'architettura più modulare.

\subsection{Ambienti di Sviluppo (IDE) e Build Automation}

Per la gestione del ciclo di vita del software sono stati utilizzati due ambienti di sviluppo integrati (IDE) complementari, scelti in base alle specifiche esigenze delle diverse fasi di lavoro:

\begin{itemize}
    \item \textbf{Visual Studio Code}: utilizzato come editor versatile per le attività di supporto e modellazione. Nello specifico, grazie all'integrazione dell'estensione \textit{UMLet}, è stato impiegato per la realizzazione dei diagrammi UML. È stato inoltre lo strumento di riferimento per la redazione della documentazione tecnica (LaTeX/Markdown) e per la gestione delle operazioni di versioning tramite l'integrazione diretta con GitHub.
    \item \textbf{Eclipse}: adottato come IDE principale per lo sviluppo sia del backend che del frontend. È stato utilizzato per la scrittura del codice sorgente Java, la gestione delle risorse web, il debugging avanzato dei microservizi e l'esecuzione dei test unitari.
\end{itemize}

La build automation e la gestione delle dipendenze del progetto Java sono state affidate ad \textbf{Apache Maven}, garantendo la riproducibilità del processo di compilazione e una gestione centralizzata delle librerie su diverse macchine di sviluppo.

\begin{figure}[H]
    \centering
    \includegraphics[width=0.6\textwidth]{2.2_toolchain.png}
    \caption{Il flusso di lavoro e la toolchain integrata del progetto JarFin.}
    \label{fig:toolchain}
\end{figure}

\subsection{Considerazioni sulla toolchain}

L’integrazione degli strumenti descritti ha permesso di coprire l’intero ciclo di vita del software, dalla modellazione iniziale alla verifica della qualità del codice.
La toolchain adottata rappresenta una soluzione completa e coerente con gli obiettivi del progetto, fornendo un supporto efficace sia allo sviluppo incrementale sia alla validazione tecnica del sistema (Figura~\ref{fig:toolchain}).
